\documentclass[11pt,a4paper]{article}

\usepackage[utf8]{inputenc}
\usepackage[T1]{fontenc}
\usepackage[brazilian]{babel}
\usepackage{lmodern}
\usepackage[a4paper,margin=2.4cm]{geometry}
\usepackage{xcolor}
\usepackage{hyperref}
\usepackage{enumitem}
\usepackage{listings}
\usepackage{parskip}
\usepackage{needspace}

\hypersetup{
  colorlinks=true,
  linkcolor=black,
  urlcolor=blue!55!black,
  pdftitle={Guia de estudo - Sprints 1 e 2},
  pdfauthor={Matheus Moreli}
}

\setlist{itemsep=3pt, topsep=5pt}
\setlength{\emergencystretch}{3em}
\lstset{
  basicstyle=\ttfamily\small,
  backgroundcolor=\color{black!4},
  frame=single,
  rulecolor=\color{black!20},
  breaklines=true,
  columns=fullflexible,
  showstringspaces=false,
  xleftmargin=4pt,
  xrightmargin=4pt
}

\newcommand{\projeto}{\paragraph{Aplicação no projeto.}}
\newcommand{\comoaplicar}{\paragraph{Como aplicar.}}
\newcommand{\referencia}[2]{\paragraph{Card e referência.}\href{#1}{#2}}
\newcommand{\cardsection}[1]{\Needspace{12\baselineskip}\subsection{#1}}

\begin{document}

\begin{titlepage}
\thispagestyle{empty}
\centering
\vspace*{3.2cm}

{\Large\textbf{SPRINTS 1 E 2}}\\[0.5cm]
{\LARGE Inventário de Ativos e Vulnerabilidades}\\[1.2cm]

{\large Material de apoio para apresentação e estudo}\\[3.2cm]

\begin{flushleft}
\large
\textbf{Aluno:} Matheus Moreli Durão Silva\\[0.35cm]
\textbf{Matrícula:} 12621CBS217\\[0.35cm]
\textbf{E-mail:} \href{mailto:matheus.moreli@ufu.br}{matheus.moreli@ufu.br}\\[0.35cm]
\textbf{Curso:} Cibersegurança - Bacharelado - Noturno
\end{flushleft}

\vfill
\end{titlepage}

\tableofcontents
\newpage

\section{Como usar este guia}

Este documento reúne as Sprints 1 e 2 do projeto de inventário de ativos e vulnerabilidades. A primeira parte apresenta os fundamentos; a segunda mostra como esses fundamentos foram reunidos no programa final. A meta não é decorar palavras de Python ou Git. A meta é conseguir responder, em cada parte do trabalho:

\begin{enumerate}
  \item O que é este conceito?
  \item Como ele funciona, passo a passo?
  \item Quando ele é útil?
  \item Onde ele aparece no programa ou no histórico do Git?
  \item Como eu demonstro que ele funciona?
\end{enumerate}

O arquivo principal do programa é \texttt{app.py}. Na Sprint 2, o código foi dividido também em \texttt{interface.py}, \texttt{inventario.py}, \texttt{modelos.py} e \texttt{persistencia.py}. Os dados persistentes ficam em \texttt{dados/inventario.json}. Este guia acompanha o repositório no GitHub e serve como material de estudo e apresentação.

\section{Visão geral do sistema}

O sistema registra ativos de TI, como notebooks, servidores, roteadores, aplicações e bancos de dados. Para cada ativo, podem ser adicionadas vulnerabilidades, isto é, problemas ou riscos de segurança encontrados naquele item.

Um ativo possui os campos ID, nome, responsável, localização, tipo e uma lista de vulnerabilidades. O ID é um número inteiro único. Ele identifica o ativo sem ambiguidade, mesmo que duas pessoas usem nomes parecidos.

\begin{lstlisting}[language=Python,caption={Formato dos dados na memória}]
ativos = {
    10: {
        "id": 10,
        "nome": "SRV-ARQUIVOS",
        "responsavel": "TI",
        "localizacao": "Sala 02",
        "tipo": 2,
        "vulnerabilidades": [
            {
                "descricao": "Sistema desatualizado",
                "categoria": "Atualizacao",
                "severidade": "Alta",
                "status": "Aberta"
            }
        ]
    }
}
\end{lstlisting}

O dicionário externo usa o ID como chave. Assim, a pergunta ``já existe um ativo com este ID?'' pode ser respondida diretamente com \texttt{if identificador in ativos}. Dentro de cada ativo, outro dicionário usa nomes claros para os campos: \texttt{ativo["responsavel"]} é mais fácil de entender do que uma posição numérica em uma lista. A lista \texttt{vulnerabilidades} é usada porque um ativo pode ter nenhuma, uma ou várias vulnerabilidades.

\section{Fluxo do programa}

Quando \texttt{executar\_programa()} começa, ela chama \texttt{carregar\_ativos()} para recuperar o JSON salvo. Em seguida, o programa permanece em um laço até a opção de saída ser escolhida.

\begin{enumerate}
  \item \textbf{Carregar:} o programa tenta ler o arquivo JSON. Caso ele não exista, começa com um dicionário vazio, \texttt{\{\}}.
  \item \textbf{Mostrar:} \texttt{mostrar\_menu()} exibe as opções no terminal.
  \item \textbf{Ler:} funções de entrada recebem a opção e os dados digitados.
  \item \textbf{Validar:} o programa verifica campos vazios, números inválidos, IDs repetidos e escolhas inexistentes.
  \item \textbf{Executar:} uma condição \texttt{if/elif/else} chama a funcionalidade escolhida.
  \item \textbf{Salvar:} depois de cadastrar, alterar ou remover, o JSON é regravado.
  \item \textbf{Sair:} a opção 0 encerra o laço com \texttt{break}.
\end{enumerate}

\section{Fundamentos de Python}

\cardsection{S1\_04 -- Saída com \texttt{print()}}

\texttt{print()} é uma função que mostra um valor no terminal. O texto entre aspas é chamado de \emph{string}. Python executa o que estiver dentro dos parênteses e envia o resultado para a tela.

\begin{lstlisting}[language=Python,caption={Formas básicas de usar print}]
# Texto simples
print("Ola, mundo!")

# Quebra de linha: \n significa nova linha
print("Inventario de ativos\n----------------------")

# Varios valores: print separa por espaco
nome = "SRV-ARQUIVOS"
identificador = 10
print("Ativo:", nome, "| ID:", identificador)

# f-string: coloca valores dentro de um texto
print(f"Ativo {nome} cadastrado com ID {identificador}.")

# Expressao dentro das chaves da f-string
print(f"Quantidade de exemplos: {2 + 3}")
\end{lstlisting}

O caractere \texttt{\textbackslash n} é uma sequência especial que cria uma nova linha. Ele ajuda a separar título e conteúdo sem precisar fazer dois \texttt{print()} diferentes. O \texttt{f} antes das aspas em uma f-string significa \emph{formatted string}: o conteúdo entre chaves é calculado e substituído pelo seu valor.

\begin{lstlisting}[language=Python,caption={Diferença entre string comum e f-string}]
nome = "Ana"
print("Ativo {nome}")   # mostra literalmente: Ativo {nome}
print(f"Ativo {nome}")  # mostra: Ativo Ana
\end{lstlisting}

\projeto
\texttt{mostrar\_menu()} utiliza \texttt{print()} para mostrar as opções. As mensagens de sucesso, de erro e de confirmação também usam essa função. Uma mensagem como ``Ativo cadastrado com sucesso'' informa o resultado da ação; sem essa saída, a pessoa não saberia se o programa executou o que foi pedido.

\referencia{https://trello.com/c/YeugQuYt}{S1\_04 no Trello}

\cardsection{S1\_05 -- Variáveis e tipos de dados}

Uma variável é um nome usado para guardar um valor durante a execução do programa. Em \texttt{identificador = 10}, o nome \texttt{identificador} passa a se referir ao número 10. O tipo do valor determina quais operações fazem sentido.

\begin{itemize}
  \item \texttt{str}: texto, como \texttt{"Sala 02"}.
  \item \texttt{int}: inteiro, como \texttt{10}.
  \item \texttt{float}: número decimal, como \texttt{3.5}.
  \item \texttt{bool}: verdadeiro ou falso, representado por \texttt{True} e \texttt{False}.
\end{itemize}

\begin{lstlisting}[language=Python,caption={Tipos de dados no contexto do projeto}]
identificador = 10          # int: ID numerico
nome = "SRV-ARQUIVOS"       # str: texto
ativo_critico = True        # bool: uma decisao sim/nao
nota = 8.5                  # float: exemplo de decimal
\end{lstlisting}

O ID foi escolhido como \texttt{int} porque o programa precisa comparar e procurar esse número. Nome, responsável e localização são \texttt{str}, pois armazenam palavras. Uma confirmação, como remover ou não remover um item, é naturalmente uma decisão booleana.

\projeto
\texttt{ler\_sim\_ou\_nao()} devolve \texttt{True} ou \texttt{False}. O valor \texttt{identificador} é um inteiro. Essa escolha deixa as condições mais claras e evita comparar textos quando o dado tem significado numérico.

\referencia{https://trello.com/c/DmvtABXt}{S1\_05 no Trello}

\cardsection{S1\_06 -- Entrada com \texttt{input()} e validação}

\texttt{input()} mostra uma pergunta e espera a pessoa digitar. Um detalhe importante: o resultado de \texttt{input()} é sempre texto. Mesmo que a pessoa digite 10, o Python recebe \texttt{"10"}. Para transformar esse texto no número inteiro 10, é preciso usar \texttt{int()}.

\begin{lstlisting}[language=Python,caption={Entrada de texto e conversão para inteiro}]
texto_id = input("Digite o ID: ")
identificador = int(texto_id)

nome = input("Digite o nome: ").strip()
if nome == "":
    print("O nome e obrigatorio.")
else:
    print(f"Voce digitou: {nome}")
\end{lstlisting}

\texttt{.strip()} remove espaços no início e no fim de um texto. Isso impede que alguém preencha um campo obrigatório apenas com espaços. A conversão \texttt{int(texto\_id)} pode gerar um erro se a pessoa digitar algo como ``dez''; esse caso é tratado na parte de exceções.

\projeto
\texttt{ler\_texto()} é reutilizada para impedir campos obrigatórios vazios. \texttt{ler\_inteiro()} tenta converter a entrada, verifica se o número é positivo e pede novamente em caso de erro. Centralizar a regra evita copiar o mesmo bloco de validação em vários lugares.

\referencia{https://trello.com/c/fFlJ9zCj}{S1\_06 no Trello}

\cardsection{S1\_07 -- Estruturas condicionais}

Uma condição permite que o programa escolha entre caminhos. \texttt{if} executa um bloco quando a expressão é verdadeira. \texttt{elif} testa outra possibilidade. \texttt{else} executa quando nenhuma condição anterior foi satisfeita.

\begin{lstlisting}[language=Python,caption={Condição para impedir ID duplicado}]
if identificador in ativos:
    print("Ja existe um ativo com esse ID.")
else:
    print("O ID pode ser usado.")
\end{lstlisting}

O operador \texttt{in} pergunta se algo pertence a uma coleção. Como o ID é uma chave do dicionário \texttt{ativos}, essa condição verifica diretamente se a chave já existe.

\projeto
O menu utiliza \texttt{if/elif/else} para chamar a função correspondente à opção. A busca decide entre encontrado e não encontrado. A remoção só ocorre após a confirmação ``SIM''. Cada situação é uma pergunta que produz verdadeiro ou falso.

\referencia{https://trello.com/c/QvZa8m7a}{S1\_07 no Trello}

\cardsection{S1\_08 -- Estruturas de repetição}

Laços de repetição evitam escrever a mesma lógica várias vezes. \texttt{while} repete enquanto uma condição for verdadeira. \texttt{for} percorre cada item de uma coleção.

\begin{lstlisting}[language=Python,caption={while e for}]
# Mantem o menu aberto ate a opcao de saida
while True:
    opcao = input("Escolha 0 para sair: ")
    if opcao == "0":
        break

# Mostra cada ativo guardado no dicionario
for ativo in ativos.values():
    print(ativo["nome"])
\end{lstlisting}

\texttt{while True} cria um laço que só termina quando o programa encontra \texttt{break}. Já \texttt{ativos.values()} entrega todos os valores do dicionário; o \texttt{for} recebe um ativo por vez na variável \texttt{ativo}.

\projeto
O laço principal mantém o menu disponível. Outros \texttt{for} mostram os tipos de ativo, percorrem a lista de vulnerabilidades e listam os ativos cadastrados.

\referencia{https://trello.com/c/op3s5sKj}{S1\_08 no Trello}

\cardsection{S1\_09 -- Funções}

Uma função é um bloco de código com nome e tarefa definida. Ela pode receber valores, chamados parâmetros, e pode devolver um resultado com \texttt{return}. Funções servem para organizar o programa e reaproveitar uma regra sem copiar e colar.

\begin{lstlisting}[language=Python,caption={Função com parâmetro e retorno}]
def nome_maiusculo(nome):
    texto_limpo = nome.strip()
    return texto_limpo.upper()

resultado = nome_maiusculo("srv-arquivos")
print(resultado)  # SRV-ARQUIVOS
\end{lstlisting}

Depois de \texttt{def}, escrevemos o nome da função, os parâmetros entre parênteses e dois-pontos. O corpo precisa ficar recuado. A variável \texttt{texto\_limpo} só é usada dentro daquela execução; o valor devolvido é guardado em \texttt{resultado}.

\projeto
\texttt{cadastrar\_ativo()}, \texttt{buscar\_ativo()}, \texttt{salvar\_ativos()} e \texttt{carregar\_ativos()} separam responsabilidades. \texttt{executar\_programa()} coordena o fluxo, sem precisar conter todos os detalhes de cadastro e arquivo.

\referencia{https://trello.com/c/XWD40Pgo}{S1\_09 no Trello}

\cardsection{S1\_10, S1\_11 e S1\_12 -- Listas, tuplas, sets e dicionários}

Coleções guardam vários valores. A estrutura deve ser escolhida pelo problema, não porque ela existe na linguagem.

\begin{lstlisting}[language=Python,caption={Comparação entre coleções}]
severidades = ["Baixa", "Media", "Alta"]
coordenadas = (10, 20)
ids_usados = {10, 20, 20}
ativo = {"id": 10, "nome": "SRV-ARQUIVOS"}

severidades.append("Critica")
print(ativo["nome"])
\end{lstlisting}

\begin{itemize}
  \item \textbf{Lista}: mantém ordem e pode mudar. \texttt{append()} adiciona ao final.
  \item \textbf{Tupla}: mantém ordem, mas não pode ser alterada depois de criada.
  \item \textbf{Set}: não guarda duplicados e não promete uma ordem fixa. No exemplo, \texttt{\{10, 20, 20\}} se torna \texttt{\{10, 20\}}.
  \item \textbf{Dicionário}: guarda pares chave--valor. \texttt{ativo["nome"]} recupera o valor usando uma chave legível.
\end{itemize}

\projeto
O inventário usa um dicionário principal com ID como chave. Cada ativo é outro dicionário com campos nomeados. O campo \texttt{vulnerabilidades} é uma lista porque novas vulnerabilidades podem ser adicionadas. Tuplas e sets não foram forçados no código porque não resolvem uma necessidade melhor que as estruturas escolhidas; saber justificar isso também demonstra entendimento.

\referencia{https://trello.com/c/EurwtLl9}{S1\_10 no Trello}; \href{https://trello.com/c/qEvT77g5}{S1\_11 no Trello}; \href{https://trello.com/c/N1J9NE4k}{S1\_12 no Trello}.

\cardsection{S1\_13 -- Enumerações}

Uma enumeração reúne opções permitidas usando nomes claros. \texttt{IntEnum} é adequada quando cada opção tem um código numérico. Isso evita usar números sem significado espalhados pelo programa.

\begin{lstlisting}[language=Python,caption={Exemplo de IntEnum}]
from enum import IntEnum

class TipoAtivo(IntEnum):
    NOTEBOOK = 1
    SERVIDOR = 2

tipo = TipoAtivo(2)
print(tipo.name)   # SERVIDOR
print(tipo.value)  # 2
\end{lstlisting}

\texttt{TipoAtivo(2)} cria o membro cujo valor é 2. \texttt{.name} mostra o nome do membro; \texttt{.value} mostra o código. Se alguém tentar \texttt{TipoAtivo(9)} e esse código não existir, Python gera \texttt{ValueError}. Esse erro é tratado para pedir uma escolha válida.

\projeto
\texttt{TipoAtivo} possui seis categorias. \texttt{escolher\_tipo\_ativo()} mostra os códigos para a pessoa escolher. O JSON salva o número, e \texttt{nome\_do\_tipo()} transforma esse código em nome legível ao exibir o ativo.

\referencia{https://trello.com/c/nJrQKljh}{S1\_13 no Trello}; \href{https://docs.python.org/3/library/enum.html}{Documentação do módulo enum}.

\cardsection{S1\_14 -- Persistência em arquivo JSON}

Dados guardados apenas em variáveis desaparecem quando o programa fecha. Persistência significa salvar esses dados em um arquivo e recuperá-los depois. JSON é um formato de texto adequado neste projeto porque representa naturalmente dicionários e listas.

\begin{lstlisting}[language=Python,caption={Gravar e ler JSON}]
import json

with open("dados/inventario.json", "w", encoding="utf-8") as arquivo:
    json.dump(ativos, arquivo, ensure_ascii=False, indent=4)

with open("dados/inventario.json", "r", encoding="utf-8") as arquivo:
    ativos = json.load(arquivo)
\end{lstlisting}

\texttt{with open(...)} abre o arquivo e o fecha automaticamente ao terminar. O modo \texttt{"w"} escreve e substitui o conteúdo anterior. O modo \texttt{"r"} lê. \texttt{json.dump()} transforma o dicionário Python em JSON; \texttt{json.load()} faz o caminho inverso. \texttt{ensure\_ascii=False} mantém acentos legíveis e \texttt{indent=4} organiza o arquivo.

Um detalhe importante: em JSON, as chaves de objetos são texto. Por isso, ao ler o arquivo, as chaves que representam IDs precisam ser convertidas de volta para inteiros antes de serem usadas pelo programa.

\projeto
\texttt{dados/inventario.json} é a base simples do trabalho. O programa salva após mudanças importantes e também ao sair. Ao iniciar novamente, os ativos e suas vulnerabilidades voltam a estar disponíveis.

\referencia{https://trello.com/c/xdRh1TRi}{S1\_14 no Trello}; \href{https://docs.python.org/3/library/json.html}{Documentação do módulo json}.

\cardsection{S1\_15 -- Tratamento de erros}

Uma exceção é um erro que interromperia o fluxo normal. \texttt{try} contém a operação que pode falhar. \texttt{except} trata um erro específico. \texttt{else} executa quando não houve erro. \texttt{raise} permite criar ou disparar uma exceção quando uma regra do sistema é violada.

\begin{lstlisting}[language=Python,caption={Tratamento de erro de conversão}]
try:
    identificador = int(input("Digite o ID: "))
except ValueError:
    print("Digite apenas um numero inteiro.")
else:
    print(f"ID valido: {identificador}")
\end{lstlisting}

Capturar um erro específico é melhor que escrever apenas \texttt{except:}. Um \texttt{except} vazio esconderia problemas diferentes e dificultaria descobrir o que aconteceu. Ao ler JSON, podem ocorrer arquivo inexistente, erro de leitura, JSON inválido ou formato inesperado. Cada situação recebe uma recuperação e uma mensagem clara.

\projeto
\texttt{carregar\_ativos()} trata \texttt{OSError}, \texttt{JSONDecodeError}, \texttt{ValueError} e \texttt{DadosInvalidosError}. A exceção personalizada \texttt{DadosInvalidosError} representa uma regra do programa: o conteúdo carregado precisa ser um dicionário de ativos. Se o arquivo estiver corrompido ou tiver formato inválido, o usuário é avisado e o programa é encerrado antes de salvar qualquer coisa. Assim, o arquivo com problema não é sobrescrito por engano.

\referencia{https://trello.com/c/qSZwgosj}{S1\_15 no Trello}; \href{https://docs.python.org/3/tutorial/errors.html}{Tutorial Python: Errors and Exceptions}.

\section{Requisitos e critérios de aceitação}

\cardsection{S1\_01 -- Compreensão do problema}

Antes de programar, é necessário entender quem usa o sistema, quais dados existem e o que precisa ser resolvido. Esse processo é chamado de compreensão de requisitos. No projeto, os elementos principais identificados foram ativo e vulnerabilidade. A pergunta principal é: como registrar os ativos de TI e acompanhar os riscos associados a cada um?

\comoaplicar
Comece listando os substantivos do problema (ativo, vulnerabilidade, responsável, localização) e depois os verbos (cadastrar, consultar, atualizar, remover). Em seguida, defina as regras que não podem ser quebradas, como manter o ID único.

\projeto
Essa análise gerou os campos guardados em cada ativo e a relação ``um ativo possui várias vulnerabilidades''. Não foi criada interface gráfica porque o objetivo da Sprint 1 é exercitar a lógica e os fundamentos de Python em um programa de terminal.

\referencia{https://trello.com/c/2nV6OG9q}{S1\_01 no Trello}

\cardsection{S1\_02 -- Requisitos funcionais}

Requisitos funcionais descrevem o que o sistema precisa fazer. Eles podem ser observados na prática. Por exemplo, ``cadastrar ativo'' significa que uma pessoa consegue informar dados e o registro passa a existir. Requisitos não funcionais descrevem qualidades, como desempenho, segurança ou facilidade de uso.

\projeto
Os verbos do enunciado se tornaram funcionalidades do menu: cadastrar, listar, buscar por ID ou nome, atualizar, remover, adicionar vulnerabilidade, visualizar vulnerabilidades e mostrar resumo. Cada opção chama uma função com tarefa definida.

\referencia{https://trello.com/c/nPnrvQad}{S1\_02 no Trello}

\cardsection{S1\_03 -- Critérios de aceitação}

Critério de aceitação é uma regra verificável que determina se uma funcionalidade está correta. Ele responde ``como saberemos que funciona?''. Um requisito sem critério pode ser entendido de formas diferentes.

\begin{itemize}
  \item Dado um ID já usado, o novo cadastro deve ser recusado e o ativo antigo deve continuar intacto.
  \item Dado um campo obrigatório vazio, o programa deve pedir novamente em vez de gravar um valor vazio.
  \item Dada uma busca sem resultado, o programa deve informar que o ativo não foi encontrado.
  \item Dada uma remoção, o programa deve pedir confirmação antes de apagar o ativo e suas vulnerabilidades.
\end{itemize}

\projeto
Esses critérios aparecem nas validações do terminal e podem ser demonstrados durante a apresentação.

\referencia{https://trello.com/c/MRXQneEZ}{S1\_03 no Trello}

\section{Git e GitHub}

\cardsection{S1\_16 -- Fundamentos de Git}

Git é um sistema de controle de versão. Ele registra a evolução do projeto localmente. GitHub é um serviço remoto que hospeda repositórios Git na internet. Um \emph{commit} é um registro de uma etapa do trabalho, acompanhado por uma mensagem que explica a mudança.

\begin{lstlisting}[caption={Sequência básica de comandos Git}]
git status
git add app.py
git commit -m "feat: cadastrar ativos"
git push
\end{lstlisting}

\begin{itemize}
  \item \texttt{git status}: mostra arquivos modificados e o estado atual.
  \item \texttt{git add}: escolhe quais alterações irão para o próximo commit.
  \item \texttt{git commit -m}: cria um registro local com uma mensagem.
  \item \texttt{git push}: envia commits para o repositório remoto.
  \item \texttt{git pull}: busca e integra alterações do remoto.
  \item \texttt{git log}: mostra a linha do tempo de commits.
\end{itemize}

\projeto
Foram criados commits separados para a aplicação, o JSON e a configuração do repositório. O repositório foi enviado ao GitHub e clonado em outro local para verificar que o remoto continha o projeto necessário.

\referencia{https://trello.com/c/vUu4H8Bq}{S1\_16 no Trello}; \href{https://git-scm.com/book/pt-br/v2}{Livro Pro Git}.

\cardsection{S1\_17 -- Versionamento e colaboração}

Uma branch é uma linha paralela de desenvolvimento. A branch \texttt{main} representa a versão principal. Uma branch de funcionalidade permite desenvolver e testar uma mudança sem alterar a versão principal imediatamente.

\begin{lstlisting}[caption={Criar e integrar uma branch}]
git switch -c feature/resumo-inventario
# alterar, testar e criar commits nesta branch
git switch main
git merge feature/resumo-inventario
\end{lstlisting}

\projeto
A funcionalidade de resumo do inventário foi desenvolvida em \texttt{feature/resumo-inventario}. Depois de testada, ela foi integrada à \texttt{main} por um merge registrado no histórico.

\referencia{https://trello.com/c/Ec5uKKhf}{S1\_17 no Trello}

\cardsection{S1\_18 -- Branches, merge e conflito}

\emph{Merge} é a integração do histórico de uma branch ao de outra. Um conflito acontece quando duas linhas de desenvolvimento alteram o mesmo trecho de modo que Git não consegue escolher sozinho qual versão usar. Conflito não é falha: é um ponto em que uma decisão humana é necessária.

Ao ocorrer um conflito, Git coloca marcadores no arquivo:

\begin{lstlisting}
<<<<<<< HEAD
texto da branch atual
=======
texto da outra branch
>>>>>>> nome-da-branch
\end{lstlisting}

Para resolver, leia as duas versões, mantenha ou combine o conteúdo correto, remova todos os marcadores, salve o arquivo e finalize o merge com um commit. Nunca deixe os marcadores no código final.

\projeto
Duas branches alteraram a mesma linha de título. O conflito foi criado para praticar, resolvido escolhendo o texto mais claro e registrado com a mensagem \texttt{merge: resolve project title conflict}.

\referencia{https://trello.com/c/kFfXpK39}{S1\_18 no Trello}

\cardsection{S1\_19 -- Pull Request e revisão de código}

Um Pull Request (PR) é uma proposta para integrar uma branch em outra no GitHub. Ele mostra a diferença entre versões, registra a explicação da mudança, permite comentários e deixa o momento da integração documentado. Em equipes, uma pessoa pode revisar a mudança feita por outra antes de aprová-la.

\comoaplicar
Ao abrir um PR, escreva o objetivo da mudança, informe como ela foi testada, reveja o diff e responda comentários. O PR deve ser integrado somente quando a alteração estiver compreensível e funcionando.

\projeto
A branch \texttt{feature/validacao-json} gerou o Pull Request \#1. A descrição explica a validação acrescentada, há um comentário de revisão e a alteração foi integrada à \texttt{main}.

\referencia{https://trello.com/c/tsLFX3pi}{S1\_19 no Trello}; \href{https://docs.github.com/pull-requests}{Documentação do GitHub sobre Pull Requests}.

\section{S1\_20 -- Revisão final e checklist}

Revisar é conseguir relacionar decisão, código e resultado. Por exemplo: ``usei um dicionário porque o ID é único e quero encontrar o ativo diretamente''. Ou: ``usei JSON porque os dados precisam continuar existindo depois que o programa fecha''.

\cardsection{Checklist de requisitos}
\begin{itemize}
  \item Explicar o problema: controlar ativos e vulnerabilidades associadas.
  \item Apontar no menu as ações de cadastrar, listar, buscar, atualizar, remover e visualizar.
  \item Demonstrar ID único, campos obrigatórios, confirmação de remoção e busca sem resultado.
\end{itemize}

\cardsection{Checklist de Python}
\begin{itemize}
  \item Executar \texttt{app.py} e navegar pelo menu.
  \item Explicar \texttt{print()}, \texttt{input()}, conversão com \texttt{int()}, condições, laços e funções.
  \item Mostrar dicionário de ativos, lista de vulnerabilidades, enumeração e arquivo JSON.
  \item Digitar uma entrada inválida e explicar como o programa recupera o erro.
\end{itemize}

\cardsection{Checklist de Git}
\begin{itemize}
  \item Abrir o histórico e explicar o que é um commit.
  \item Mostrar \texttt{main} e as branches de funcionalidade.
  \item Localizar o merge e explicar o conflito resolvido.
  \item Abrir o Pull Request \#1 e explicar o seu ciclo: branch, mudança, revisão e integração.
\end{itemize}

\cardsection{Sequência para a apresentação}
\begin{enumerate}
  \item Apresente o problema e os dados de um ativo.
  \item Cadastre um ativo e uma vulnerabilidade.
  \item Consulte o ativo e mostre a vulnerabilidade vinculada.
  \item Mostre uma validação, como ID duplicado ou entrada não numérica.
  \item Feche e abra o programa para provar que o JSON preserva os dados.
  \item Abra o GitHub e mostre commits, branches, merge e Pull Request.
\end{enumerate}

\referencia{https://trello.com/c/X3LzSoi0}{S1\_20 no Trello}

\clearpage
\section{Sprint 2 -- Construção do projeto final}

Na Sprint 1, os conceitos foram praticados em exercícios pequenos: entrada, saída, condições, listas, dicionários, JSON e Git. A Sprint 2 junta essas peças em um programa único. O foco deixa de ser ``como usar um comando isolado'' e passa a ser ``como organizar uma funcionalidade completa sem tornar o código complicado''.

O programa continua sendo de terminal e usa apenas recursos básicos de Python: funções, classes simples, \texttt{Enum}, dicionários, listas, arquivo JSON e \texttt{try/except}. Não há banco de dados, interface gráfica, framework ou consulta automática à internet. Essa escolha é intencional: o problema é pequeno e precisa ser fácil de explicar.

\cardsection{Organização simples do código}

O projeto foi separado por responsabilidade. Separar não significa criar muitas camadas: significa evitar que um mesmo arquivo faça tudo ao mesmo tempo.

\begin{itemize}
  \item \texttt{app.py}: inicia o programa e informa o caminho do JSON.
  \item \texttt{interface.py}: contém \texttt{input()}, \texttt{print()} e o menu.
  \item \texttt{inventario.py}: contém regras de cadastro, consulta, atualização e exclusão.
  \item \texttt{modelos.py}: define \texttt{Ativo}, \texttt{Vulnerabilidade} e \texttt{TipoAtivo}.
  \item \texttt{persistencia.py}: lê e grava o arquivo JSON.
\end{itemize}

\begin{lstlisting}[language=Python,caption={Fluxo principal simplificado}]
ativos = carregar_ativos("dados/inventario.json")
executar_programa(ativos, "dados/inventario.json")
\end{lstlisting}

O menu não precisa conhecer os detalhes de JSON, e a função que salva JSON não precisa fazer perguntas ao usuário. Quando cada parte tem uma tarefa principal, fica mais fácil localizar erros e explicar o código.

\cardsection{S2\_01 e S2\_02 -- Boas práticas, classes e Enum}

Uma classe é um molde para representar dados relacionados. Neste trabalho, \texttt{Ativo} reúne ID, nome, responsável, localização, tipo, criticidade, vulnerabilidades e histórico. \texttt{Vulnerabilidade} reúne CVE, CWE, CVSS, descrição, fonte, data, impacto, prioridade, tratamento, status e verificação.

\begin{lstlisting}[language=Python,caption={Ideia de uma classe simples}]
class Ativo:
    def __init__(self, identificador, nome):
        self.id = identificador
        self.nome = nome
        self.vulnerabilidades = []
\end{lstlisting}

\texttt{self} representa o objeto que está sendo criado ou alterado. Assim, dois ativos podem existir ao mesmo tempo, cada um com seus próprios dados. Uma função como \texttt{cadastrar\_ativo()} recebe os dados, verifica regras e cria um objeto. Ela não usa \texttt{input()} nem \texttt{print()}; isso fica na interface.

\texttt{TipoAtivo} é um \texttt{IntEnum}. Ele associa nomes claros a códigos inteiros, como \texttt{SERVIDOR = 2}. Isso atende ao requisito de possuir categorias padronizadas sem espalhar números sem significado pelo código.

\referencia{https://trello.com/c/ALQSY6p9}{S2\_01 no Trello}; \href{https://trello.com/c/hu3YSI6L}{S2\_02 no Trello}.

\cardsection{S2\_03, S2\_08 a S2\_11 -- CRUD de ativos}

CRUD é uma sigla para quatro operações sobre dados: \emph{Create} (criar), \emph{Read} (consultar), \emph{Update} (atualizar) e \emph{Delete} (excluir). O dicionário \texttt{ativos} usa o ID como chave. Isso permite localizar um ativo diretamente.

\begin{lstlisting}[language=Python,caption={Busca direta pela chave do dicionário}]
def consultar_por_id(ativos, identificador):
    return ativos.get(identificador)
\end{lstlisting}

\texttt{.get()} devolve o valor encontrado ou \texttt{None} quando a chave não existe. Isso evita erro ao buscar um ID ausente. No cadastro, o programa rejeita ID duplicado. Na atualização, o ID não é modificável: nome, responsável, localização, tipo e criticidade podem mudar, mas a identidade do ativo é preservada. Na exclusão, a tela mostra o alvo e exige a palavra \texttt{SIM}; ao remover o ativo, sua lista de vulnerabilidades também desaparece junto.

\referencia{https://trello.com/c/NdgkJaYg}{S2\_03 no Trello}.

\cardsection{S2\_12 a S2\_14 -- Vulnerabilidades}

Uma vulnerabilidade é uma fragilidade; ameaça é algo que pode explorar uma fragilidade; risco combina a possibilidade do problema com seu impacto. Por isso, o programa não usa apenas uma nota CVSS: ele também registra impacto no ativo, prioridade, tratamento, status e como a correção será verificada.

Cada vulnerabilidade pertence a um ativo. O par ``ID do ativo + CVE'' a identifica de forma simples: o mesmo CVE pode afetar ativos diferentes, mas não pode aparecer duas vezes no mesmo ativo. As quatro operações também existem para vulnerabilidades: cadastrar, consultar, atualizar e remover.

\begin{lstlisting}[language=Python,caption={Validação simples do CVSS}]
nota = float(cvss)
if not 0.0 <= nota <= 10.0:
    raise DadosInvalidosError("CVSS deve estar entre 0.0 e 10.0.")
\end{lstlisting}

O \texttt{raise} interrompe apenas a operação inválida e a interface mostra uma mensagem compreensível. O menu continua aberto. O sistema também exige uma verificação antes de mudar o status para ``Corrigida'', pois dizer que algo foi corrigido sem evidência não é suficiente.

Como referência de consulta externa, o registro \href{https://nvd.nist.gov/vuln/detail/CVE-2021-44228}{CVE-2021-44228 no NVD} descreve uma vulnerabilidade do Apache Log4j com CVSS 10.0. O programa não consulta o NVD automaticamente: o usuário registra manualmente fonte e data, mantendo a aplicação simples e funcionando sem internet.

\referencia{https://trello.com/c/zzyFLROU}{S2\_12 no Trello}; \href{https://trello.com/c/xMcbgbvn}{S2\_13 no Trello}; \href{https://trello.com/c/XFF3yzHz}{S2\_14 no Trello}.

\cardsection{S2\_04 a S2\_07 -- Requisitos e validações}

Um requisito diz o que o sistema deve fazer. Um critério de aceite diz como comprovar que ele fez corretamente. Por exemplo, ``cadastrar ativo'' é o requisito; ``o ID precisa ser inteiro, positivo e único, e o ativo deve aparecer na busca'' é o critério verificável.

\begin{itemize}
  \item Campo obrigatório vazio: o programa pede novamente, sem fechar.
  \item ID repetido: o ativo antigo não é alterado e o novo não é cadastrado.
  \item Busca inexistente: aparece ``Ativo não encontrado'', sem erro técnico.
  \item Exclusão cancelada: nada é removido.
  \item JSON inválido: o programa não sobrescreve o arquivo corrompido com uma base vazia.
\end{itemize}

Essas regras aparecem em funções de validação, como \texttt{validar\_texto()}, \texttt{validar\_cve()}, \texttt{validar\_cwe()} e \texttt{validar\_cvss()}. Quando uma regra falha, elas usam \texttt{DadosInvalidosError}; a interface captura essa exceção e apresenta uma mensagem curta.

\referencia{https://trello.com/c/NcSDMJuv}{S2\_04 no Trello}; \href{https://trello.com/c/YSC57Mjy}{S2\_05 no Trello}; \href{https://trello.com/c/5beR1TIT}{S2\_06 no Trello}; \href{https://trello.com/c/Vs045ALu}{S2\_07 no Trello}.

\cardsection{S2\_15 -- Dicionários como índice}

O dicionário é a estrutura central do inventário:

\begin{lstlisting}[language=Python,caption={Estrutura principal na memória}]
ativos = {
    10: ativo_servidor,
    20: ativo_notebook
}
\end{lstlisting}

Em vez de percorrer todos os itens para achar o ID 10, o código usa \texttt{ativos.get(10)}. A chave é única, o formato dos objetos é uniforme e as alterações preservam a mesma chave. Ao salvar, os objetos viram dicionários simples; ao carregar, voltam a ser objetos \texttt{Ativo} e \texttt{Vulnerabilidade}.

\cardsection{S2\_18 -- Testes e correções}

\texttt{test\_sprint2.py} não é um segundo programa para o usuário usar. É uma coleção de verificações automáticas feita com o módulo padrão \texttt{unittest}, que já vem com Python. Cada teste cria uma base pequena na memória, chama uma função do inventário e compara o resultado esperado com o resultado obtido.

\begin{lstlisting}[caption={Como executar os testes}]
python -m unittest -v test_sprint2.py
\end{lstlisting}

Os 14 testes verificam cadastro, busca, ID duplicado, tipo inválido, filtros, atualização e confirmação de exclusão de ativos. Também verificam CVE duplicado, CVE e CVSS inválidos, atualização e remoção de vulnerabilidades, exigência de verificação para correção e proteção do JSON corrompido. Os testes usam arquivos temporários e não alteram \texttt{dados/inventario.json}.

Além dos testes automáticos, vale demonstrar o fluxo real no terminal: cadastrar um ativo, vincular uma vulnerabilidade, consultar, atualizar e remover. Assim, a avaliação observa tanto as regras isoladas quanto o programa em uso.

\cardsection{S2\_16, S2\_17 e S2\_19 -- GitHub e apresentação}

O código está no GitHub, na branch \texttt{main}. O histórico possui branches de funcionalidade, commits com mensagens explicativas e merges por Pull Request. A tag \texttt{trabalho-1-entrega} identifica uma versão preparada para apresentação.

Para demonstrar em cinco minutos:
\begin{enumerate}
  \item Execute \texttt{python app.py}.
  \item Cadastre um ativo e uma vulnerabilidade associada.
  \item Consulte o ativo e mostre que a vulnerabilidade está vinculada a ele.
  \item Atualize o responsável ou o status da vulnerabilidade.
  \item Tente repetir um ID ou informar CVSS inválido para mostrar a validação.
  \item Remova uma vulnerabilidade, primeiro cancelando e depois confirmando.
  \item Abra o GitHub, mostre \texttt{main}, as Pull Requests e a tag final.
\end{enumerate}

Ao explicar, prefira frases concretas: ``usei um dicionário porque o ID é único e a busca fica direta''; ``usei JSON porque os dados continuam depois que o programa fecha''; ``separei interface e regras para não misturar perguntas com validação''.

\referencia{https://trello.com/c/UvtQ6Mc6}{S2\_16 no Trello}; \href{https://trello.com/c/FBSYXNJP}{S2\_17 no Trello}; \href{https://trello.com/c/ADSVtVQ4}{S2\_19 no Trello}.

\section{Referências gerais das Sprints 1 e 2}

\begin{itemize}
  \item Cards S1\_01 a S1\_20 e S2\_01 a S2\_19 do quadro Trello da disciplina.
  \item Pressman e Maxim, \emph{Engenharia de Software}.
  \item Sommerville, \emph{Engenharia de Software}.
  \item Eric Matthes, \emph{Curso Intensivo de Python}.
  \item \href{https://docs.python.org/3/}{Documentação oficial do Python}.
  \item \href{https://docs.python.org/3/library/unittest.html}{Documentação do módulo unittest}.
  \item \href{https://git-scm.com/book/pt-br/v2}{Pro Git}.
  \item \href{https://docs.github.com/}{Documentação do GitHub}.
  \item \href{https://nvd.nist.gov/vuln/detail/CVE-2021-44228}{NVD -- CVE-2021-44228}.
\end{itemize}

\end{document}
